TEAM MEMBERS

TEAMS (1 to 3 people)

PRESENTATIONS:

Presentations will be held in the last week of classes, Apr 20th and Apr 22nd.

Duration: approximately 5-10mins per team (depends on the number of teams).

Each team member must present for approximately equal amount of time.

FINAL REPORT

Instructions on the final report. Each group only needs to submit one report by one of your team members in PDF format. 

1. The due time will be 11:59pm, Sun, May3rd via CANVAS. (only one team member needs to submit)   
Late reports are accepted up to 24hr late, with 1% penalty applied for each hour of delay (maximum: 24% penalty).
2. You should follow the NeurIPS style for your final report. Here you can find the template and latex source file:   https://nips.cc/Conferences/2018/PaperInformation/StyleFiles Links to an external site.Links to an external site..
Submit in PDF format.
3. There is no minimum limit on the number of pages, but the maximum is 8 pages following the NeurIPS style. You can submit a supplementary file, but the TA is not obligated to read through your supplementary file. Therefore, the main report should contain all essential parts of your work. Upload your code as a separate file (R/Python/Julia).
The report should include the following:
1. Purpose of the project. Description of the problem you are studying and the purpose of and rationale of the project.
- descriptive statistics and informative exploratory data analysis of the dataset
- description of any data cleaning that's been applied to the dataset
- research question/questions
2. Literature Review. Review of the literature related to the proposed project. The literature should address the analysis done on similar data using frequentist and/or Bayesian approaches. 
3. Proposed Method. A written description of the project plan including
* the research question,
* the Bayesian model
* details of the computational solution (MCMC details). 
(should include some if not all out of the following)
  - sample size for the MCMC/MC simulation
  - priors and parameters of the prior distributions
  - proposal distribution and its parameter/parameters
  - effective sample size, ACFs, posteriors for relevant parameters, posterior CIs/means/SDs
You are not allowed to use the method outside of the scope of the class (No U Turn sampling, Hamiltonian Monte Carlo etc) make sure to describe the method.

4. Data analysis and results interpretation. 

When you submit your final report, please attach your R/Python/Julia code (well documented so that we can easily re-run your code). 
Please feel free let me know if you have any further questions. 